\chapter{The Thermal Process}

\section{Installation Overview}

The installation studied in this project is a laboratory-scale thermal rig located at
ICAM Engineering School. It consists of three main sub-systems: an electric furnace,
a plate heat exchanger, and two fluid circuits equipped with flow measurement and
proportional valve control.

\textbf{Electric furnace.}
The furnace is an electrically-heated recirculation unit. It heats the hot fluid circuit
(F1) and circulates it at a regulated flow rate. The furnace is controlled via a Boolean
output signal (\plcaddr{Q0.3}). Its outlet temperature, denoted
\code{T\_hout} (equivalently \code{Tin1\_HE1} at the exchanger inlet), is the primary
energy input to the heat exchanger. The return temperature \code{T\_hin} represents the
hot-circuit fluid after passing through the exchanger.

\textbf{Plate heat exchanger HE-001.}
The heat exchanger is a gasketed plate-type unit. Hot fluid (F1) enters at $T_{in1,HE1}$,
flows through the hot channels, and exits at $T_{out1,HE1}$. Cold fluid (F2) enters at
$T_{in2,HE1}$ from the opposite end and exits at $T_{out2,HE1}$, in a counter-current
configuration. This arrangement maximises the temperature driving force and thermal
efficiency.

\textbf{Fluid circuits.}
The hot circuit~F1 is a closed loop between the furnace and the hot side of HE-001.
The cold circuit~F2 is a separate loop with independent flow control.
Each circuit is equipped with an electromagnetic flow meter (outputs \code{F1} and
\code{F2}, in \si{\metre\cubed\per\hour}) and a proportional valve controlled by an
analog output from the PLC (\code{Valve\_F1} at \plcaddr{QW8}; \code{Valve\_F2} at
\plcaddr{QW4}).

Figure~\ref{fig:pid_diagram} shows a simplified P\&ID of the installation.

\begin{figure}[H]
  \centering
  \begin{tikzpicture}[
    box/.style={rectangle, draw=darkblue, thick, fill=lightgray!30,
                minimum width=2.6cm, minimum height=1.2cm, font=\small, align=center},
    sensor/.style={circle, draw=darkblue, thick, fill=white,
                   minimum size=0.7cm, font=\tiny},
    flow/.style={-Stealth, thick, draw=orange!80!black},
    coldflow/.style={-Stealth, thick, draw=blue!70},
    arr/.style={-Stealth, thick, draw=darkblue},
  ]
    % Furnace
    \node[box, fill=orange!15] (furnace) at (0,0)
      {\textbf{Furnace}\\power (\plcaddr{Q0.3})};

    % Heat exchanger
    \node[box, fill=blue!8, minimum width=3.2cm, minimum height=2.4cm]
      (HE) at (6,0) {\textbf{HE-001}\\plate\\exchanger};

    % Hot circuit (top)
    \draw[flow] (furnace.east) -- ++(1.0,0)
      node[above, font=\scriptsize]{$T_{in1}$~\plcaddr{MD48}} -- (HE.west|-HE.north)
      +(0,-0.4);

    \draw[flow] (HE.east|-HE.north) +(0,-0.4) -- ++(1.2,0)
      node[above, font=\scriptsize]{$T_{out1}$~\plcaddr{MD44}} -- ++(0.5,0);

    % Return hot (bottom of hot path)
    \draw[flow] ($(furnace.east)+(2.2,0)$) to[out=270,in=270,looseness=0.5]
      node[below, font=\scriptsize]{$T_{hin}$~\plcaddr{MD16}} (furnace.south|-furnace.east)
      (furnace.south);

    % Cold circuit (bottom) — counter-current
    \draw[coldflow] (HE.east|-HE.south) +(0,0.4) -- ++(2.0,0)
      node[below, font=\scriptsize]{$T_{out2}$~\plcaddr{MD52}} ;

    \draw[coldflow] ($(HE.west)+(0,-0.4)$) -- ++(-2.0,0)
      node[below, font=\scriptsize]{$T_{in2}$~\plcaddr{MD40}};

    % Flow sensors
    \node[sensor] (F1s) at (3.8, 0.55) {F1};
    \node[sensor] (F2s) at (3.8,-0.55) {F2};

    % Labels
    \node[font=\scriptsize, above=0.1 of furnace] {Hot circuit F1};
    \node[font=\scriptsize, below=0.8 of HE]
      {Cold circuit F2 (counter-current)};
  \end{tikzpicture}
  \caption{Simplified P\&ID of the laboratory thermal installation.}
  \label{fig:pid_diagram}
\end{figure}

\section{Instrumentation and Signals}

\subsection{Temperature Sensors}

Temperatures are measured by industrial sensors (thermocouple or PT100 resistance
thermometer). The sensor output is a 4--20\,mA current signal read by an analog input
module of the S7-1500. The module converts this signal to a 16-bit integer in the range
0--27\,648, stored in a \plcaddr{IW} address (Process Image Input Word). A scaling
function in OB1 then converts the raw integer to a temperature value in \si{\degreeCelsius}
using the linear formula:
\begin{equation}
  T \;[\si{\degreeCelsius}]
  = \frac{\text{raw}}{27\,648} \times (T_{\max} - T_{\min}) + T_{\min}.
  \label{eq:scaling}
\end{equation}
The result is stored in a \plcaddr{MD} address (Merker Double-word, 32-bit IEEE~754 float).

\subsection{Flow Meters}

The electromagnetic flow meters output a 4--20\,mA signal proportional to volumetric flow
rate. The same scaling formula~\eqref{eq:scaling} converts the raw \plcaddr{IW} integer to
\si{\metre\cubed\per\hour}. The hot circuit flow $F_1$ is stored at \plcaddr{MD28}
(raw from \plcaddr{IW316}), and the cold circuit flow $F_2$ at \plcaddr{MD32}
(raw from \plcaddr{IW204}).

\subsection{Valve Actuators}

The proportional valves accept an analog command (0--20\,mA or 0--10\,V). The PLC writes
a Word value (0--27\,648, corresponding to 0--100\,\% opening) to the analog output address:
\code{Valve\_F1} at \plcaddr{QW8} and \code{Valve\_F2} at \plcaddr{QW4}.
The OB1 program converts the engineering-unit flow setpoints ($F_{1,SP}$ in
\si{\metre\cubed\per\hour}, stored at \plcaddr{MD24}) to the appropriate Word output via
the inverse of~\eqref{eq:scaling}.

\subsection{Complete Process Tag Table}

Table~\ref{tab:tags} lists the 21 process tags retained after reviewing the full tag table
(25 intermediate or redundant variables were removed). The naming convention uses English
names with the suffix \code{\_HE1} for heat exchanger variables.

\begin{longtable}{@{}L{4.6cm}L{2.2cm}L{1.5cm}L{1.5cm}L{3cm}@{}}
  \caption{Complete WinCC Unified tag table (21 process tags).}
  \label{tab:tags}\\
  \toprule
  \textbf{Tag name} & \textbf{PLC addr.} & \textbf{Type} & \textbf{Unit} & \textbf{Description} \\
  \midrule
  \endfirsthead
  \multicolumn{5}{l}{\small\itshape (continued)}\\
  \toprule
  \textbf{Tag name} & \textbf{PLC addr.} & \textbf{Type} & \textbf{Unit} & \textbf{Description} \\
  \midrule
  \endhead
  \bottomrule
  \endfoot
  \code{T\_hin}                      & \plcaddr{MD16}  & Real & \si{\degreeCelsius} & Furnace inlet temperature \\
  \code{T\_hout}                     & \plcaddr{MD12}  & Real & \si{\degreeCelsius} & Furnace outlet temperature \\
  \code{T\_we\_rock\_furnace}        & \plcaddr{IW300} & Int  & ADC     & Furnace inlet raw signal \\
  \code{T\_wy\_Oven\_scale}          & \plcaddr{IW304} & Int  & ADC     & Furnace outlet raw signal \\
  \code{Tin1\_HE1}                   & \plcaddr{MD48}  & Real & \si{\degreeCelsius} & HE hot inlet temperature \\
  \code{Tout1\_HE1}                  & \plcaddr{MD44}  & Real & \si{\degreeCelsius} & HE hot outlet temperature \\
  \code{Tin2\_HE1}                   & \plcaddr{MD40}  & Real & \si{\degreeCelsius} & HE cold inlet temperature \\
  \code{Tout2\_HE1}                  & \plcaddr{MD52}  & Real & \si{\degreeCelsius} & HE cold outlet temperature \\
  \code{Tin1\_HE1\_raw}              & \plcaddr{IW112} & Int  & ADC     & $T_{in1}$ raw ADC \\
  \code{Tout1\_HE1\_raw}             & \plcaddr{IW120} & Int  & ADC     & $T_{out1}$ raw ADC \\
  \code{Tin2\_HE1\_raw}              & \plcaddr{IW116} & Int  & ADC     & $T_{in2}$ raw ADC \\
  \code{Tout2\_HE1\_raw}             & \plcaddr{IW108} & Int  & ADC     & $T_{out2}$ raw ADC \\
  \code{F1}                          & \plcaddr{MD28}  & Real & \si{\metre\cubed\per\hour} & Hot flow rate \\
  \code{F1\_SP}                      & \plcaddr{MD24}  & Real & \si{\metre\cubed\per\hour} & Hot flow setpoint \\
  \code{F1\_skal}                    & \plcaddr{IW316} & Int  & ADC     & $F_1$ raw ADC \\
  \code{Valve\_F1}                   & \plcaddr{QW8}   & Word & 0--27648 & Hot valve command \\
  \code{F2}                          & \plcaddr{MD32}  & Real & \si{\metre\cubed\per\hour} & Cold flow rate \\
  \code{F2\_SP}                      & \plcaddr{MD36}  & Real & \si{\metre\cubed\per\hour} & Cold flow setpoint \\
  \code{F2\_skal}                    & \plcaddr{IW204} & Int  & ADC     & $F_2$ raw ADC \\
  \code{Valve\_F2}                   & \plcaddr{QW4}   & Int  & 0--27648 & Cold valve command \\
  \code{power}                       & \plcaddr{Q0.3}  & Bool & ---     & Furnace power (ON=1) \\
\end{longtable}

\section{Process Physics}

\subsection{Heat Exchanger Thermal Model}

Energy transfer in a counter-current plate heat exchanger is governed by Newton's law of
cooling applied across the heat transfer surface:
\begin{equation}
  Q = U \cdot A \cdot \text{LMTD},
  \label{eq:heat_transfer}
\end{equation}
where $Q$ [\si{\watt}] is the total heat transfer rate, $U$ [\si{\watt\per\metre\squared\per\kelvin}]
is the overall heat transfer coefficient, $A$ [\si{\metre\squared}] is the total heat
transfer area, and LMTD [\si{\kelvin}] is the Log Mean Temperature Difference.

For a counter-current configuration:
\begin{equation}
  \text{LMTD} = \frac{\Delta T_1 - \Delta T_2}{\ln(\Delta T_1 / \Delta T_2)}
  \quad \text{with} \quad
  \Delta T_1 = T_{in1} - T_{out2},\quad
  \Delta T_2 = T_{out1} - T_{in2}.
  \label{eq:lmtd}
\end{equation}

At steady state, the energy balance on the cold side gives:
\begin{equation}
  Q = \dot{m}_2 \cdot C_{p2} \cdot (T_{out2} - T_{in2}).
  \label{eq:cold_balance}
\end{equation}
This shows that $T_{out2}$ depends on both the heat transfer rate $Q$ (driven by $T_{in1}$)
and the cold flow rate $F_2$. An increase in $T_{in1}$ raises LMTD, increases $Q$, and
therefore raises $T_{out2}$.

\subsection{Dynamic Behaviour and Cascade Motivation}

From a control perspective, the dynamic path from the actuator ($F_{1,SP}$) to the primary
controlled variable ($T_{out2,HE1}$) is sequential:

\begin{enumerate}[leftmargin=*, itemsep=3pt]
  \item \textbf{Inner dynamic} ($F_{1,SP} \rightarrow T_{in1,HE1}$):
    Increasing $F_{1,SP}$ opens Valve\_F1, increasing hot flow rate. More fluid circulates
    through the furnace per unit time. $T_{in1}$ responds with a first-order-like rise
    characterised by time constant $\tau_1$ and dead time $\theta_1$.

  \item \textbf{Outer dynamic} ($T_{in1,HE1} \rightarrow T_{out2,HE1}$):
    The hotter fluid entering the exchanger increases LMTD, increases $Q$, and heats
    the cold fluid more strongly. $T_{out2}$ rises with its own time constant $\tau_2 > \tau_1$
    (the exchanger integrates thermal energy over its plate surface area).
\end{enumerate}

The combined transfer function from $F_{1,SP}$ to $T_{out2}$ is approximately the product
of two first-order systems, producing a higher-order response with a large apparent dead
time. A single PID closing the loop around this combined dynamic must be tuned
conservatively, resulting in sluggish disturbance rejection.

\textit{Cascade control} exploits the measurability of $T_{in1,HE1}$. By regulating $T_{in1}$
with a fast inner PI loop, disturbances in the hot circuit are corrected before they reach
$T_{out2}$. The outer controller handles only the slower HE dynamics, operating at a
much lower bandwidth. This structure is the primary motivation for the cascade strategy
adopted in this project.

\subsection{Linearised FOPTD Model}

For controller design purposes, the inner process ($F_{1,SP} \rightarrow T_{in1,HE1}$) is
linearised around a nominal operating point and approximated by a First-Order Plus Time Delay
(FOPTD) transfer function:
\begin{equation}
  G_{\text{inner}}(s) = \frac{K \cdot e^{-\theta s}}{\tau s + 1},
  \label{eq:foptd}
\end{equation}
where $K$ [\si{\degreeCelsius\per(\metre\cubed\per\hour)}] is the static gain, $\tau$
[\si{\second}] the time constant, and $\theta$ [\si{\second}] the dead time.
This approximation is valid in the neighbourhood of the operating point and for step
amplitudes not exceeding 20\,\% of the operating range.

\section{Control Objectives}

The control system must satisfy the following requirements:

\paragraph{Primary objective.}
Regulate $T_{out2,HE1}$ to a setpoint $T_{out2,SP}$ (e.g., \SI{45}{\degreeCelsius})
with acceptable steady-state error and settling time.

\paragraph{Disturbance rejection.}
Attenuate the effect of furnace temperature variations (changes in $T_{hin}$) and cold
flow rate changes ($F_2$ variations) before they propagate to $T_{out2}$.

\paragraph{Actuator constraints.}
$F_{1,SP}$ must remain within the safe operating range $[F_{1,\min},\, F_{1,\max}]$.
Valve commands must be smooth to avoid mechanical wear from rapid oscillations.

\paragraph{Operating modes.}
The system supports three modes of operation:
MANUAL (operator sets $F_{1,SP}$ directly),
AUTO-INNER (inner PI+P loop active, operator sets $T_{in1,SP}$), and
AUTO-FULL (full cascade active, operator sets $T_{out2,SP}$).

\section{Nominal Operating Point}

\pendingbox{%
  This section will be completed with numerical values measured during the experimental
  session of 2026-05-19 with Michal.%
}

\begin{table}[H]
  \centering
  \caption{Nominal operating conditions (values to be measured).}
  \label{tab:oppoint}
  \begin{tabular}{@{}L{4.5cm}C{2.5cm}C{2cm}L{4cm}@{}}
    \toprule
    \textbf{Variable} & \textbf{Nominal value} & \textbf{Unit} & \textbf{Notes} \\
    \midrule
    $T_{hin}$       & TBD & \si{\degreeCelsius} & Steady-state furnace return \\
    $T_{in1,HE1}$   & TBD & \si{\degreeCelsius} & Steady-state hot inlet \\
    $T_{out2,HE1}$  & TBD & \si{\degreeCelsius} & Target cold outlet \\
    $F_1$           & TBD & \si{\metre\cubed\per\hour} & Hot flow at op.\ point \\
    $F_2$           & TBD & \si{\metre\cubed\per\hour} & Cold flow (held constant) \\
    \bottomrule
  \end{tabular}
\end{table}
